\section{Primera Entrevista}
\begin{itemize}
  \item ¿Se desea controlar el legajo de alumnos y docentes
  solamente?
  \item ¿Sería ideal tener una tabla para ver las condiciones?
  \item ¿Podría enlazarse cada requisito del legajo (su archivo
  en la base de datos) de cada docente/alumno en los campos de
  la tabla?
  \item ¿Hay información extra necesaria de cada perfil? ¿Cuál?
  \item Si hubieran notificaciones sobre el legajo, ¿cuál sería
  el medio ideal para eso?
  \item ¿Cuáles son los formularios que podrán hacer los docentes
  y los alumnos?
  \item ¿Qué información es necesaria sobre los exalumnos?
  \item ¿Por qué se necesita documentación histórica de los
  exalumnos?
  \item ¿Cómo se rellenan los formularios?
  \item ¿Cuáles son las acciones que tendrán permitidas cada uno
  de los usuarios?
  \item ¿Qué partes del legajo pueden digitalizarse?
  \item ¿Es preferible manejar un sistema de notificaciones o
  mensajería directa?
\end{itemize}
\section{Segunda Entrevista}
\begin{itemize}
  \item ¿Cómo se clasifican los docentes, según su salario?
  \item ¿Cuál es la cantidad de horas máximas que puede trabajar
  un docente?
  \item ¿Quién autoriza el rol de un docente?
  \item ¿Qué tipos de \glspl{reldoc} existen?
  \item ¿Qué diferencias hay entre los tipos de relevos?
  \item Un docente nuevo que releva a otro por renuncia a las
  horas, ¿puede necesitar un cambio de horarios?
  \item ¿De quién depende el manejo de los horarios?
  \item ¿Cuáles horarios pueden rotar dentro del año?
  \item ¿Qué condiciones deben cumplirse para el cambio de
  horarios?
  \item ¿Cómo manejan los cambios de horario?
  \item ¿De qué forma verifican actualmente el cumplimiento de
  las condiciones para cambios de horario?
  \item ¿Qué razones podrían haber para que el sistema pueda
  validar cambios de horario?
\end{itemize}
